\documentclass[10pt,a4paper]{article}

\usepackage[utf8]{inputenc}
\usepackage[T1]{fontenc}
\usepackage[spanish,english]{babel}
\usepackage{amsfonts,amsmath,graphicx,array,tabularx}

\newcommand{\titlesize}{\fontsize{18}{20pt}\bfseries}
\newcommand{\titulo}[1]{\vspace*{10pt}\begin{center}
               \titlesize{#1}\par\vspace*{10pt}\large}
\newcommand{\autor}[1]{\vspace*{12pt}\sl{#1}\par}
\def\affilnum#1{${}^{#1}$}
\def\affil#1{${}^{#1}$}
\newcommand{\address}[1]{\vspace*{10pt}\footnotesize
                \textrm{#1}\end{center}\par\vspace*{20pt}}

\newcommand{\Acknowledgements}[1]{\noindent {\bf {\normalsize Acknowledgements:}}
{\normalsize #1}  }

%%%%%%%%%%%%%%%%%%%%%%%%%%%%%%%%%%%%%%%%%%%%%%%%%%%%%%%%%%%%%%%%%%%%%%

\begin{document}

% Title of the talk
\titulo
{A Machine-Learning Approach to the %Attainable 
Jordan Block Structure %Under Bounded Similarity Conditioning
}

% Names of the authors. Underline the author who gives the talk
\autor
{\underline{Michał Trojanowski}\affil{1}, Michał Wojtylak\affil{2}}

% Affiliation of authors
\address
{\affilnum{1}\ Faculty of Mathematics and Computer Science, Jagiellonian University, Poland\\
E-mail: {\tt michal.trojanowski@student.uj.edu.pl\qquad} \\
 \affilnum{2}\ Faculty of Mathematics and Computer Science, Jagiellonian University, Poland\\
E-mail: {\tt michal.wojtylak@uj.edu.pl\qquad} }


%% Abstract
\begin{center}
{\bf \large Abstract}
\end{center}

The size of the largest Jordan block associated with a clustered eigenvalue plays a significant role in applied mathematics. Classical perturbation theory \cite{wilkinson1965}, \cite{Demmel1983}, as well as more recent works on coalescence theory \cite{Alam2011}, \cite{Banks2021}, establish that large Jordan blocks occur only in  proximity of eigenvector ill-conditioning. % In particular, large algebraic multiplicities require high similarity condition numbers.

Motivated by this constraint, we study the following problem: given a matrix whose eigenvalues form a cluster and whose eigenvector matrix satisfies a prescribed conditioning bound, determine the maximal Jordan block size attainable under perturbations of controlled magnitude. Following Demmel's dissociation framework, we restrict attention to a single eigenvalue cluster and construct matrices of the form $ A = S^{-1} (J+\Delta) S$, where $J$ has arbitrary Jordan structure for a single eigenvalue, $\Delta$ is a perturbation of controlled size, and $\kappa(S)$ is bounded.

We present a transformer-based model trained on such a dataset of matrices, with core components that are dimension-independent. The model predicts the maximal block size with accuracy exceeding 95\% for small perturbations in dimension 35. Moreover, we observe decreasing accuracy as the perturbation magnitude increases, indicating a transition from geometrically structured spectral behaviour to a noise-dominated regime.

Our results provide a data-driven complement to classical distance-to-defectivity theory and reveal a conditioning-induced regularity in near-defective matrices.

\bigskip

\noindent {\bf \large References}

\renewcommand{\refname}{}
\vspace*{-1.25cm}

\begin{thebibliography}{5}

\bibitem{wilkinson1965}
J. H. Wilkinson. The Algebraic Eigenvalue Problem. Monographs on Numerical Analysis. Oxford
University Press, Oxford, UK, 1965.

\bibitem{Demmel1983}
J. W. Demmel. \textit{A Numerical Analyst’s Jordan Canonical Form}. PhD thesis, University of California, Berkeley, 1983.

\bibitem{Alam2011}
 R. Alam, S. Bora, R. Byers, and M. L. Overton. Characterization and construction of the nearest
defective matrix via coalescence of pseudospectral components. Linear Algebra and its Applications,
435(3):494–513, 2011.

\bibitem{Banks2021}
J. Banks, A. Kulkarni, S. Mukherjee, and N. Srivastava. Gaussian regularization of the pseudospectrum
and davies’ conjecture. Communications on Pure and Applied Mathematics, 74(10):2114–2131, 2021.

\end{thebibliography}

\end{document}
